\documentclass[conference,compsoc]{IEEEtran}

\usepackage{amsmath,amssymb,amsthm}
\usepackage{graphicx}
\usepackage{booktabs}
\usepackage{hyperref}
\usepackage{xcolor}
\usepackage{multirow}
\usepackage{microtype}

% ── Theorem environments ──────────────────────────────────────────────────────
\newtheorem{theorem}{Theorem}
\newtheorem{corollary}{Corollary}
\newtheorem{proposition}{Proposition}
\newtheorem{definition}{Definition}
\newtheorem{lemma}{Lemma}

% ── TODO marker (remove before submission) ────────────────────────────────────
\newcommand{\TODO}[1]{\textcolor{red}{\textbf{[TODO: #1]}}}

% ── Paper metadata ────────────────────────────────────────────────────────────
\begin{document}

\title{What Makes a Dataset Privacy-Sensitive?\\
  Towards Predicting Membership Inference Risk Before Model Training}

\author{
  \IEEEauthorblockN{Anonymous Submission}
  \IEEEauthorblockA{IEEE S\&P 2027 Submission \#TODO}
}

\maketitle

% =============================================================================
\begin{abstract}
% Write abstract LAST — after all results are in.
% Key elements: 1 sentence problem, 3 findings with numbers, 1 sentence implication.
%
Membership inference attacks (MIA) have been studied for a decade, yet we
still lack a principled answer to a more fundamental question: why do some
datasets leak more than others, independent of the model trained on them?
We introduce the \emph{Dataset Privacy Risk Index} (DPRI), a
pre-training risk estimator built from five data-distributional statistics
that require no trained model.
Across \TODO{7} datasets spanning four domains, DPRI predicts MIA
vulnerability with $R^2 = $\TODO{X.XX} and Spearman $\rho = $\TODO{X.XX}
under leave-one-dataset-out cross-validation.
Our analysis yields three security findings:
(1)~dataset structure explains \TODO{XX}\% of MIA AUC variance, versus
\TODO{XX}\% for model architecture;
(2)~widely used MIA benchmarks (Purchase100, Texas100) are structural
outliers that overstate real-world privacy risk by \TODO{X.X}$\times$ on
average; and
(3)~the benefit of differential privacy depends primarily on DPRI, not
on the nominal $\varepsilon$ value, implying that current DP calibration
practices may systematically under-protect high-risk datasets.
We release DPRI as an open-source tool for pre-training dataset auditing.
\end{abstract}

% =============================================================================
% =============================================================================
\section{Introduction}
\label{sec:intro}
% =============================================================================

Over the past decade, membership inference attacks
(MIA)~\cite{homer2008resolving,shokri2017membership} have become the
standard tool for measuring how much a trained model leaks about its
training data.
Hundreds of papers have proposed stronger
attacks~\cite{salem2019ml,nasr2019comprehensive,choquette2021label,song2021systematic,carlini2022lira,ye2022enhanced,watson2022importance,zarifzadeh2024rmia},
dedicated defenses~\cite{nasr2018machine,jia2019memguard,abadi2016deep},
tighter theoretical
bounds~\cite{yeom2018privacy,sablayrolles2019whitebox}, and
systematizations of the resulting
landscape~\cite{papernot2018sok,hu2022membership}.
The same model-centric view organizes the sibling lines of work:
attribute and property
inference~\cite{fredrikson2015model,ganju2018property}, training-data
extraction~\cite{carlini2019secret,carlini2021extracting}, differential
privacy~\cite{dwork2006calibrating,abadi2016deep}, and machine
unlearning~\cite{cao2015towards,bourtoule2021machine,chen2021unlearning}
all quantify or remove leakage from a trained model.
Yet a more fundamental question has gone largely unanswered:

\smallskip
\begin{center}
Why do some datasets leak more than others,\\
independent of the model trained on them?
\end{center}
\smallskip

This is not only an engineering question.
It is also a question about the nature of data itself.
The pre-ML privacy literature established that re-identification risk is
driven by intrinsic properties of the data: the sparsity of the Netflix
ratings matrix~\cite{narayanan2008robust}, the uniqueness of mobility
traces~\cite{demontjoye2013unique} and credit-card
metadata~\cite{demontjoye2015unique}, and the incompleteness that
defeats naive anonymization~\cite{sweeney2002k,rocher2019estimating}.
Machine learning privacy has largely set this lesson aside, treating the
dataset as a fixed backdrop behind the model. We argue that membership
inference is, in part, a symptom of how the training data is shaped.

\paragraph{The problem with current practice.}
Every MIA paper chooses a dataset, trains a model, and reports an AUC.
The AUC is then attributed to the attack method or the defense mechanism.
But when we decompose the variation in AUC across 31 datasets, three
attacks, and three model families, the dataset is anything but a fixed
backdrop. Among the regularized models that practitioners actually deploy,
dataset choice explains $49.5\%$ of the variation in AUC, twice the share
of model architecture ($24.4\%$) and well above the attack method
($26.1\%$). Only a deliberately overfit model, an unbounded-depth Random
Forest that memorizes its training set outright, erases this margin
(Section~\ref{sec:finding1}).
Dataset structure is a first-class driver of privacy risk,
yet it is rarely reported or controlled for.

A second pathology is benchmark reliance.
Purchase100 and Texas100~\cite{shokri2017membership} are the two most
common MIA benchmarks.
We show that both are geometric outliers. Their sample uniqueness is the
highest among all 31 datasets we study, placing them far outside the range
of typical deployed data, so papers that report results only on these
benchmarks measure privacy risk in an unrepresentatively high-risk regime.

\paragraph{Our approach.}
We ask a temporally prior question:
Can the privacy risk of a dataset be anticipated before any model is trained?
If so, practitioners gain an actionable signal: estimate risk at data
collection time, before committing to a training pipeline or a privacy budget.

We give a qualified yes. We introduce the Dataset Privacy Risk Index (DPRI),
a scalar pre-training risk estimator derived from distributional statistics:
sample uniqueness, local density, outlier score, feature entropy, class
cluster separation, and data dimensionality.
Each statistic is motivated by a theoretical result
(Theorem~\ref{thm:main}) showing that MIA advantage is bounded by a
function of local data geometry.
Across 31 datasets, DPRI ranks datasets by risk with a Spearman correlation
of $0.61$ ($p < 0.001$). We report this number under nested cross-validation
that selects features on training folds only, so it is free of the selection
optimism that inflates small-$n$ results; the same discipline is what
exposed our own early estimates, on fewer datasets, as artifacts.
Data geometry significantly shapes privacy risk without fully determining
it.

Our question is related to, but distinct from, the memorization literature.
Feldman~\cite{feldman2020memorization} showed that rare individual samples
must be memorized for good generalization, a phenomenon observable only
after training. We ask the complementary and temporally prior question:
can dataset-level structural statistics predict aggregate membership
inference risk before any model is trained? Where that line of work
identifies which samples a given model will memorize (post-training,
sample-level), we predict how much a dataset will leak (pre-training,
dataset-level), a distinction that is both scientific and operational. A
pre-training score is the only kind a practitioner can act on at data
collection time.

\paragraph{Contributions.}
\begin{itemize}

\item \textbf{Theoretical foundation.}
We prove that the Bayes-optimal MIA advantage for a sample $x$ is
upper-bounded by $O(u(x) / \rho(x)^{1/d})$, where $u(x)$ is sample
uniqueness and $\rho(x)$ is local density, both computable from raw data
(Theorem~\ref{thm:main}).
This provides the first formal link between pre-training data geometry
and membership inference risk.

\item \textbf{Dataset Privacy Risk Index (DPRI).}
We define a dataset-level risk score from six distributional features and
validate it across 31 datasets in five domains using three MIA attacks and
three model families. DPRI ranks datasets by risk with an unbiased nested
cross-validated Spearman correlation of $0.61$ ($p < 0.001$).

\item \textbf{A calibrated, qualified finding.}
The prediction is moderate, not precise, and no single statistic suffices.
Sample uniqueness, local density, class separability, and dimensionality
are each significant predictors, yet privacy risk is shaped by data geometry
without being determined by it. We argue this nuance is itself a
contribution. It bounds what any pre-training predictor can achieve.

\item \textbf{Two measurement findings.}
(F1)~Among realistic, regularized models, dataset structure governs MIA
risk more than model architecture ($49.5\%$ vs.\ $24.4\%$ of AUC variation)
and far more than the choice of attack ($26.1\%$), so an attack number
reported on a single benchmark conflates three distinct sources of
variation, and the dataset is the largest of them.
(F2)~The benchmarks that anchor a decade of MIA research are geometric
outliers, casting doubt on the external validity of results reported only
on them.

\item \textbf{Open-source tool.}
We release DPRI as a Python package that audits any tabular dataset in
minutes on a laptop, with no model training and no GPU, outputting a risk
score and its distributional breakdown.

\end{itemize}

\paragraph{Scope.}
We study 31 public datasets spanning tabular, image, network, medical, and
recommendation domains, in the centralized training setting.
Extensions to federated learning and generative models are left as future
work.
All datasets used are publicly available. No proprietary data was accessed.

% =============================================================================
% Section 2: Background and Threat Model
% (sec_threat_model.tex content merged here)
% =============================================================================

\section{Background and Threat Model}
\label{sec:background}

\subsection{Membership Inference Attacks}

Let $D = \{(x_i, y_i)\}_{i=1}^{n}$ be a training dataset drawn i.i.d.\
from an unknown distribution $\mathcal{P}$ over $\mathcal{X} \times \mathcal{Y}$.
A learner trains a classifier $f_\theta : \mathcal{X} \to \Delta(\mathcal{Y})$
by minimizing empirical risk on $D$.
A membership inference attack (MIA) is an algorithm
$\mathcal{A}$ that, given a target sample $(x, y)$ and query access to
$f_\theta$, outputs $1$ if $(x,y) \in D$ and $0$ otherwise.

The MIA advantage of $\mathcal{A}$ is:
\begin{equation}
  \mathrm{Adv}(\mathcal{A}, f_\theta, D)
  =
  \bigl|
    \Pr[\mathcal{A}=1 \mid (x,y)\in D]
    -
    \Pr[\mathcal{A}=1 \mid (x,y)\notin D]
  \bigr|.
  \label{eq:adv}
\end{equation}
We denote the dataset-level risk as the expected advantage over samples:
\begin{equation}
  \mathrm{Risk}(D) = \sup_{\mathcal{A}} \;\mathbb{E}_{(x,y)\sim\mathcal{P}}\,
  \mathrm{Adv}(\mathcal{A}, f_\theta, D).
  \label{eq:risk}
\end{equation}
In practice, $\mathrm{Risk}(D)$ is estimated via AUC of the best
empirical attack over a grid of attacks and model families
(Section~\ref{sec:setup}).

\subsection{Differential Privacy}

An algorithm $\mathcal{M}$ is $(\varepsilon,\delta)$-differentially
private~\cite{dwork2006calibrating} if for all adjacent datasets $D, D'$
differing in one record and all outputs $S$:
$\Pr[\mathcal{M}(D)\in S] \leq e^\varepsilon \Pr[\mathcal{M}(D')\in S] + \delta$.
Yeom et al.~\cite{yeom2018privacy} showed that
$\mathrm{Adv}(\mathcal{A}^*) \leq e^\varepsilon - 1 + \delta$
for $(\varepsilon,\delta)$-DP mechanisms.
Our work shows that this bound, while tight in the worst case, does not
account for how much DPRI amplifies or attenuates the practical risk
(Section~\ref{sec:finding3}).

\subsection{Threat Model}
\label{sec:threat}

We consider two principals: a data holder who trains a model on
$D$ and may later publish it, and an adversary who infers membership.

\paragraph{Adversary goal.}
Given $(x,y)$ and oracle access to $f_\theta$, determine whether
$(x,y) \in D$.

\paragraph{Adversary capability.}
Black-box access: the adversary submits inputs and observes
predicted class probabilities, but cannot inspect model weights.
This is the standard setting for deployed APIs and MLaaS platforms.
White-box results are consistent and reported in Appendix~\ref{app:whitebox}.

\paragraph{Adversary knowledge.}
The adversary knows the data domain distribution $\mathcal{P}$
(e.g., that the dataset contains medical records of a specific demographic)
but does not know the exact set $D$.

\paragraph{Defender goal.}
Before training, compute $\mathrm{DPRI}(D)$ from raw data alone, with no
model required, to estimate $\mathrm{Risk}(D)$ and inform DP budget
selection and release decisions.

\paragraph{Dual use.}
DPRI also has an offensive application: an adversary with access to
candidate target datasets can compute DPRI pre-attack and prioritize
high-risk datasets to maximize attack return on query budget.
We quantify this in Section~\ref{sec:implications} and discuss mitigations.

\paragraph{Scope.}
Centralized training, passive adversary (no data poisoning).
Extensions to federated learning are discussed in Section~\ref{sec:implications}.

% ============================================================
% Section 3: Theoretical Analysis
% ============================================================

\section{Theoretical Analysis}
\label{sec:theory}

We establish that MIA advantage is upper-bounded by data-distributional
quantities computable before training.
All proofs are in Appendix~\ref{app:proofs}.

\subsection{Notation}

Let $D = \{x_1,\dots,x_n\} \subset \mathcal{X}$ be drawn i.i.d.\ from
density $p$ on $\mathcal{X} \subseteq \mathbb{R}^d$.
For any $x \in D$, define:
\begin{align}
  u(x) &= \min_{x' \in D,\, x'\neq x} \|x - x'\|_2
           \quad \text{(sample uniqueness / nearest-neighbour distance)},
           \label{eq:uniqueness}\\
  \rho(x) &= \frac{k}{ n \cdot \mathrm{Vol}(B(x,\, r_k(x)))}
           \quad \text{(local density, $k$-NN estimate)},
           \label{eq:density}
\end{align}
where $r_k(x)$ is the distance to the $k$-th nearest neighbour of $x$ in
$D \setminus \{x\}$ and $\mathrm{Vol}(\cdot)$ denotes Euclidean ball volume.

For the in-distribution and out-distribution of a trained model
$f_\theta$, define the score distributions:
\begin{equation}
  P_{\mathrm{in}}(s) = \Pr[\ell(f_\theta, x) \leq s \mid x \in D], \quad
  P_{\mathrm{out}}(s) = \Pr[\ell(f_\theta, x) \leq s \mid x \notin D],
\end{equation}
where $\ell(f_\theta,x)$ is the model's loss on sample $x$.

\subsection{Main Theorem}

\begin{theorem}[DPRI Upper Bound on MIA Advantage]
\label{thm:main}
Let $f_\theta$ be trained on $D$ by minimising empirical risk.
Assume the loss function $\ell$ is $L$-Lipschitz in $x$.
Then for any sample $x \in D$:
\begin{equation}
  \mathrm{Adv}(\mathcal{A}^*, f_\theta, x)
  \;\leq\;
  \sqrt{\frac{L^2 \cdot u(x)^2}{2\,\sigma^2}}
  \;\cdot\;
  \frac{1}{\rho(x)^{1/d}},
  \label{eq:bound}
\end{equation}
where $\mathcal{A}^*$ is the Bayes-optimal adversary,
$\sigma^2$ is the variance of the loss under $P_{\mathrm{out}}$,
and $d$ is the data dimensionality.

At the dataset level, the aggregate advantage satisfies:
\begin{equation}
  \mathrm{Risk}(D)
  \;\leq\;
  \frac{1}{n}\sum_{x \in D}
  \sqrt{\frac{L^2 \cdot u(x)^2}{2\,\sigma^2\,\rho(x)^{2/d}}}.
  \label{eq:risk_bound}
\end{equation}
\end{theorem}

\paragraph{Interpretation.}
The bound says that privacy risk increases when samples are
far from their neighbours ($u(x)$ large, i.e.\ isolated)
and when the local population is sparse ($\rho(x)$ small).
Both quantities are computable directly from $D$, with no model training required.
The Lipschitz constant $L$ and noise level $\sigma$ are model-family
constants that can be estimated or treated as fixed hyperparameters.

\begin{proof}[Proof sketch]
We proceed in three steps.

Step 1: KL to advantage.
By the Bretagnolle--Huber inequality, the Bayes-optimal MIA advantage is:
\begin{equation}
  \mathrm{Adv}(\mathcal{A}^*) \leq
  \sqrt{\tfrac{1}{2}\,\mathrm{KL}(P_{\mathrm{in}} \| P_{\mathrm{out}})}.
  \label{eq:bh}
\end{equation}

Step 2: KL in terms of loss gap.
Since $P_{\mathrm{in}}$ and $P_{\mathrm{out}}$ are loss score distributions,
and $\ell$ is $L$-Lipschitz, the expected loss gap between a training
sample $x$ and the nearest out-of-sample point $x'$ satisfies:
\begin{equation}
  \mathbb{E}[\ell(f_\theta, x)] - \mathbb{E}[\ell(f_\theta, x')]
  \;\leq\; L \cdot \|x - x'\|_2
  \;\leq\; L \cdot u(x).
  \label{eq:lipschitz_gap}
\end{equation}
Under a sub-Gaussian model for $P_{\mathrm{out}}$ with variance $\sigma^2$:
\begin{equation}
  \mathrm{KL}(P_{\mathrm{in}} \| P_{\mathrm{out}})
  \;\leq\;
  \frac{(L \cdot u(x))^2}{2\,\sigma^2}.
  \label{eq:kl_bound}
\end{equation}

Step 3: Density correction.
The nearest-neighbour distance $u(x)$ and local density $\rho(x)$ are
related by the $k$-NN geometry: for a sample in a region of density
$\rho$, the expected distance to the nearest neighbour scales as
$\rho(x)^{-1/d}$ (standard result, see \cite{biau2015lectures}).
Substituting into~\eqref{eq:kl_bound} and applying~\eqref{eq:bh}
yields~\eqref{eq:bound}.
Averaging over all $x \in D$ gives~\eqref{eq:risk_bound}.
\qed
\end{proof}

\subsection{Discussion of the Bound}

\paragraph{Tightness.}
The bound is tight up to constants for datasets with a uniform distribution
over a compact manifold (Proposition~\ref{prop:tight}, Appendix~\ref{app:proofs}).
For real-world datasets, the bound overestimates risk by a factor that
decreases as the data distribution becomes more regular.
Section~\ref{sec:results} shows empirically that the bound ranks datasets
in the same order as the observed AUC (Spearman $\rho > 0.85$), suggesting
it is tight enough to be predictively useful.

\paragraph{What the bound cannot capture.}
The bound depends on $L$ (Lipschitz constant of the loss) and $\sigma$
(out-distribution noise), which are model-family properties.
This means the bound explains why dataset structure is the dominant
driver of risk, while model architecture contributes a fixed multiplicative
constant, consistent with our empirical Finding~1 (Section~\ref{sec:finding1}).

\paragraph{Relationship to differential privacy.}
If the training mechanism is $(\epsilon,\delta)$-DP, then
$\mathrm{Adv}(\mathcal{A}^*) \leq e^\epsilon - 1 + \delta$~\cite{yeom2018privacy}.
Combined with Theorem~\ref{thm:main}, high-DPRI datasets require
smaller $\epsilon$ to achieve the same risk reduction, motivating our
data-adaptive DP calibration recommendation (Section~\ref{sec:finding3}).

\subsection{Corollary: Benchmark Bias is Predictable}

\begin{corollary}[Benchmark datasets have structurally elevated DPRI]
\label{cor:benchmark}
If a benchmark dataset $D_b$ was constructed by subsampling or clustering
a larger dataset to create balanced, well-separated classes, then
$\mathrm{DPRI}(D_b) > \mathrm{DPRI}(D_{\mathrm{real}})$
for a real-world dataset $D_{\mathrm{real}}$ drawn from the same domain,
in expectation over the construction procedure.
\end{corollary}

\begin{proof}[Proof sketch]
Subsampling increases $u(x)$ (nearest neighbours are removed) and
clustering increases class separation, both of which increase the RHS
of~\eqref{eq:risk_bound}.
\end{proof}

This corollary provides the theoretical underpinning for our empirical
Finding~2 (benchmark bias, Section~\ref{sec:finding2}).

% =============================================================================
\section{Dataset Privacy Risk Index}
\label{sec:dpri}
% =============================================================================

Motivated by Theorem~\ref{thm:main}, we define five distributional
features that are computable from raw data before any model training.
Each feature targets a geometric property that the theorem identifies as
a driver of MIA advantage.

\subsection{Feature Definitions}

Let $D = \{x_1,\dots,x_n\}$ denote the feature matrix (labels used only
for Feature~5).
All per-sample scores are aggregated to dataset level using mean and
90th percentile, yielding a feature vector
$\phi(D) \in \mathbb{R}^{8}$ (two statistics for each of the four
per-sample features, plus two scalars).

\paragraph{Feature 1: Sample Uniqueness $u(x)$.}
\begin{equation}
  u(x) = \min_{x' \in D,\, x'\neq x} \|x - x'\|_2
\end{equation}
High $u(x)$ means $x$ is isolated from all other training points.
Isolated samples have high MIA advantage by Equation~\eqref{eq:bound}.
We use the 5-NN distance as a robust estimate. The nearest neighbour
alone ($k=1$) is noisy for dense datasets.

\paragraph{Feature 2: Local Density $\rho(x)$.}
\begin{equation}
  \rho(x) = \frac{1}{r_k(x)}, \quad
  r_k(x) = \text{distance to } k\text{-th nearest neighbour}
\end{equation}
Low $\rho(x)$ (sparse region) implies large $u(x)$ and high risk.
The two features are theoretically related but empirically complementary:
$u(x)$ captures global isolation while $\rho(x)$ captures local structure.

\paragraph{Feature 3: Outlier Score $o(x)$.}
\begin{equation}
  o(x) = \frac{1}{2}\bigl(\mathrm{LOF}_{\mathrm{norm}}(x)
                         + \mathrm{iForest}_{\mathrm{norm}}(x)\bigr)
\end{equation}
An ensemble of Local Outlier Factor~\cite{breunig2000lof} and Isolation
Forest~\cite{liu2008isolation}, each normalized to $[0,1]$.
The ensemble reduces the sensitivity of any single outlier detector to
its specific density assumption.

\paragraph{Feature 4: Feature Entropy $H(X)$.}
\begin{equation}
  H(X) = \frac{1}{d} \sum_{j=1}^{d} H(X_j), \quad
  H(X_j) = -\sum_b p_b \log_2 p_b
\end{equation}
The average Shannon entropy across all features $X_j$, computed over
a histogram discretization.
High-entropy features increase the dimensionality of the space in which
samples are embedded, reducing density and increasing uniqueness.

\paragraph{Feature 5: Cluster Separation $S$.}
\begin{equation}
  S = \mathrm{Silhouette}(X, y)
  = \frac{1}{n}\sum_{i=1}^{n} \frac{b(x_i) - a(x_i)}{\max(a(x_i), b(x_i))}
\end{equation}
where $a(x_i)$ is the mean intra-class distance and $b(x_i)$ is the
mean nearest-class distance.
Well-separated classes ($S$ near $1$) place decision boundary samples
in high-contrast regions, increasing their MIA risk.

\subsection{Aggregation and the DPRI Score}

The per-sample scores $(u, \rho, o)$ are aggregated to dataset level as
mean and 90th percentile, giving a feature vector:
\begin{equation}
  \phi(D) = \bigl(
    \bar{u},\; u_{90},\;
    \bar{\rho},\; \rho_{90},\;
    \bar{o},\; o_{90},\;
    H(X),\; S
  \bigr) \in \mathbb{R}^{8}.
\end{equation}
The 90th percentile captures tail behaviour (the hardest-to-protect
samples) which disproportionately drives aggregate MIA AUC.

The final DPRI score is a linear combination:
\begin{equation}
  \mathrm{DPRI}(D) = w^\top \phi(D),
  \label{eq:dpri}
\end{equation}
where $w \in \mathbb{R}^{8}$ is learned by ridge regression to predict
$\mathrm{Risk}(D)$ (Section~\ref{sec:setup}).
The linear form is intentional: it is interpretable, and
Section~\ref{sec:results} shows it performs comparably to nonlinear
regressors, indicating the relationship between DPRI features and
$\mathrm{Risk}(D)$ is approximately linear.

\subsection{Computation Cost}

All five features require only standard nearest-neighbour and
density-estimation routines.
On a laptop (8-core CPU), DPRI computation takes \TODO{X} seconds for
datasets up to 50k samples and \TODO{X} minutes for 200k samples.
No GPU is required.
This is the key practical advantage over MIA-based risk assessment,
which requires training and attacking multiple models.

% =============================================================================
\section{Experimental Setup}
\label{sec:setup}
% =============================================================================

\subsection{Datasets}

We evaluate on 31 publicly available datasets spanning five domains: medical,
image, network, tabular, and recommendation or mobility. They range from
several hundred to roughly 200k samples and from 5 to 6169 features. The set
includes the two canonical MIA benchmarks, Purchase100 and Texas100, which
we include specifically to test Finding~2. Datasets larger than 30k samples
are subsampled to 30k to keep the attack-model grid tractable. Most are
loaded through scikit-learn~\cite{pedregosa2011scikit},
OpenML~\cite{vanschoren2014openml}, or the UCI
repository~\cite{dua2019uci}. Table~\ref{tab:all_datasets} lists the
full corpus with sources, and Table~\ref{tab:datasets} details a
representative selection.

\begin{table}[t]
\centering
\caption{The full 31-dataset corpus. OpenML datasets are identified by
         their stable numeric \texttt{data\_id}; the remainder come from
         scikit-learn loaders or their original public releases.}
\label{tab:all_datasets}
\small
\setlength{\tabcolsep}{2.2pt}
\begin{tabular}{lll}
\toprule
\textbf{Dataset} & \textbf{Domain} & \textbf{Source} \\
\midrule
Texas100      & hospital benchmark   & Shokri et al.~\cite{shokri2017membership} \\
Purchase100   & retail benchmark     & Shokri et al.~\cite{shokri2017membership} \\
Adult         & census income        & UCI~\cite{adult1996} \\
COMPAS        & criminal justice     & ProPublica~\cite{compas2016} \\
Pima Diabetes & clinical             & UCI~\cite{smith1988pima} \\
Breast-W      & clinical             & OpenML 15 \\
JM1, KC1      & software defects     & OpenML 1053, 1067 \\
MovieLens-1M  & recommendation       & GroupLens~\cite{movielens} \\
Gowalla       & mobility             & SNAP~\cite{cho2011friendship} \\
MNIST         & image                & OpenML 554~\cite{lecun1998gradient} \\
Fashion-MNIST & image                & OpenML 40996 \\
Digits        & image                & scikit-learn \\
OptDigits     & image                & OpenML 28 \\
PenDigits     & image                & OpenML 32 \\
Letter        & image                & OpenML 6 \\
Segment       & image segmentation   & OpenML 36 \\
SatImage      & remote sensing       & OpenML 182 \\
CovType       & forest cover         & scikit-learn~\cite{blackard1999comparative} \\
German Credit & credit               & OpenML 31~\cite{dua2019uci} \\
Bank Marketing& marketing            & OpenML 1461 \\
Spambase      & e-mail               & OpenML 44 \\
Nomao         & web entities         & OpenML 1486 \\
Electricity   & energy               & OpenML 151 \\
Mushroom      & categorical          & OpenML 24 \\
Magic         & gamma telescope      & OpenML 1120 \\
Phoneme       & speech               & OpenML 1489 \\
HAR           & motion sensors       & OpenML 1478 \\
Gas Drift     & chemical sensors     & OpenML 1476 \\
Vehicle       & silhouettes          & OpenML 54 \\
Ionosphere    & radar                & OpenML 59 \\
\bottomrule
\end{tabular}
\end{table}

The corpus size is a deliberate design choice. Because the dataset is our
unit of statistical analysis, every dataset added to the corpus requires
its own full ground-truth grid of nine attack-model configurations,
including shadow-model training for LiRA, so each point in our regression
costs hours of compute. Thirty-one datasets is roughly an order of
magnitude more than the two to four typically used in MIA evaluations,
and Section~\ref{sec:power} shows empirically why that breadth is not
optional: correlations measured on small corpora can be perfect and
spurious at once.

\begin{table}[t]
\centering
\caption{Representative datasets (of 31). Sample and feature counts are the
         original sizes after preprocessing; datasets larger than 30k
         samples are subsampled to 30k for the attack grid.}
\label{tab:datasets}
\small
\setlength{\tabcolsep}{2pt}
\begin{tabular}{llrrr}
\toprule
\textbf{Dataset} & \textbf{Domain} & \textbf{Samples} & \textbf{Feat.} & \textbf{Cls.} \\
\midrule
Texas100~\cite{shokri2017membership}    & Benchmark/med.  & 67,330 & 6,169 & 100 \\
Purchase100~\cite{shokri2017membership} & Benchmark/retail & 197,324 & 600 & 100 \\
MNIST~\cite{lecun1998gradient}  & Image       & 30,000 & 784 & 10 \\
CovType~\cite{blackard1999comparative} & Forest & 30,000 &  54 &  7 \\
Adult~\cite{adult1996}          & Income      & 48,842 & 14 & 2 \\
German Credit~\cite{dua2019uci} & Credit      &  1,000 & 74 & 2 \\
Spambase~\cite{dua2019uci}      & Spam        &  4,601 & 57 & 2 \\
Gowalla~\cite{cho2011friendship}& Mobility    & 107,092 &  6 & 2 \\
\bottomrule
\end{tabular}
\end{table}

\subsection{Preprocessing}

All datasets pass through one deterministic pipeline. Categorical
features and labels are integer-encoded, every feature is standardized
to zero mean and unit variance, and datasets above 30k samples are
subsampled once, with a fixed seed, before any further processing.
Standardization matters twice over: it is conventional for the models,
and it puts the geometric features of Section~\ref{sec:dpri} on a
common scale so that nearest-neighbour distances are comparable across
datasets. Each dataset is then split 50\%/50\% into members and
non-members by a stratified split with a fixed seed, the standard MIA
evaluation convention; the target model trains on the member half, and
attack AUC is computed over the full member and non-member sets.

\subsection{Ground Truth: Risk(D)}

For each dataset we run three MIA attacks against three model families,
yielding nine AUC scores, and define $\mathrm{Risk}(D)$ as their mean:
\begin{equation}
  \mathrm{Risk}(D) = \frac{1}{|\mathbb{A}||\mathbb{F}|}
  \sum_{\mathcal{A} \in \mathbb{A}} \sum_{f \in \mathbb{F}}
  \mathrm{AUC}(\mathcal{A}, f, D),
  \label{eq:risk_gt}
\end{equation}
where $\mathbb{A} = \{\text{LiRA, shadow, loss-threshold}\}$ and
$\mathbb{F} = \{\text{MLP, XGBoost, RF}\}$.
We use the mean rather than the maximum because the maximum collapses onto
the Random Forest for every dataset (Section~\ref{sec:finding1}), which
saturates near AUC $1.0$ and erases the contrast between datasets. The mean
reflects the typical adversary across realistic and worst-case models alike.

A second metric choice deserves justification. Carlini et
al.~\cite{carlini2022lira} argue that individual attacks should be
judged by true-positive rate at low false-positive rate, because an
attacker who confidently identifies a few members matters more than one
who is weakly right on average. That argument concerns evaluating
\emph{attacks}; our target is a property of \emph{datasets}. We need a
per-dataset scalar that is stable across nine heterogeneous
attack-model configurations and comparable across datasets whose member
sets range from a few hundred to fifteen thousand samples. AUC
integrates over all thresholds, is invariant to the 50/50 class balance
we fix by construction, and remains estimable at our smallest dataset
sizes, where a TPR at $0.1\%$ FPR rests on a handful of samples and is
dominated by noise. Because DPRI predicts a ranking, any monotone risk
functional would do in principle; mean AUC is the one that is stable
enough at $n = 31$ datasets to regress against.

The three attacks span the spectrum from strongest to simplest.
LiRA~\cite{carlini2022lira}, the current state of the art, uses 16
shadow models to compute a per-sample likelihood ratio; the shadow-model
attack~\cite{shokri2017membership} trains 4 shadow models and learns a
binary meta-classifier; and the loss-threshold
attack~\cite{yeom2018privacy} thresholds the model's confidence on the
correct class, which is Bayes-optimal under a Gaussian assumption. The
three model families are an MLP (two hidden layers, 128-64, trained with
Adam~\cite{kingma2015adam} for 100 epochs on GPU),
XGBoost~\cite{chen2016xgboost} (300 trees, depth 6), and a Random
Forest~\cite{breiman2001random} (300 trees, unbounded depth). All models
use a 50\%/50\% member/non-member split, with seeds fixed at 42.
The grid is chosen for coverage rather than exhaustiveness: the attacks
span the simplest baseline to the published state of the art
(Section~\ref{sec:attack_families}), and the model families span the
three regimes that dominate tabular practice, a regularized neural
network, a regularized boosted ensemble, and a deliberately
unregularized memorizer that serves as the adversarial worst case.

\subsection{DPRI Computation}

We compute the five geometric features using scikit-learn and scipy, with
$k=5$ for the $k$-NN features, and append data dimensionality $\log d$. No
model training is required. Silhouette score uses a 5,000-point subsample for
datasets with $n > 5{,}000$.

\subsection{Regression and Validation}

We predict $\mathrm{Risk}(D)$ from the DPRI features under
leave-one-dataset-out cross-validation across all 31 datasets. Features are
rank-transformed within each fold (Section~\ref{sec:dpri}). We report Spearman
rank correlation as the primary metric, together with an unbiased nested
cross-validation that selects the feature subset on the training folds only.
We also report $R^2$, but it is not the appropriate lens here, since the
features are rank-transformed and the prediction target is a ranking.

\subsection{Reproducibility}

All code and data-processing scripts are provided as an anonymized
artifact and will be released publicly upon publication.
All random seeds are fixed, and results are deterministic given the same
hardware.

% =============================================================================
\section{Results}
\label{sec:results}
% =============================================================================

We evaluate on 31 datasets across five domains, each attacked with three MIA
methods (LiRA, shadow, loss-threshold) under three model families (MLP,
XGBoost, Random Forest), for 279 (dataset, attack, model) configurations.

\subsection{Experimental Ground Truth}

We define the ground-truth risk of a dataset as the mean MIA AUC over the
nine attack-model configurations, $\mathrm{Risk}(D)$.
We use the mean rather than the per-dataset maximum because the maximum
collapses onto the Random Forest for every dataset. Random Forest with no
depth limit memorizes its training set, so its AUC saturates near $1.0$ and
erases all contrast between datasets. The mean preserves that contrast.

Aggregated over the grid, LiRA is the strongest attack (mean AUC $0.71$)
ahead of loss-threshold and shadow ($0.61$ each), and Random Forest is the
most vulnerable model (mean AUC $0.75$) ahead of XGBoost ($0.62$) and MLP
($0.56$). $\mathrm{Risk}(D)$ ranges almost two-fold across datasets, from
$0.92$ (Texas100) down to $0.51$ (Gowalla, Mushroom), confirming that the
choice of dataset alone moves privacy risk substantially.

\subsection{DPRI Ranks Privacy Risk}

We predict $\mathrm{Risk}(D)$ from the DPRI features using ridge regression
under leave-one-dataset-out cross-validation. Because the geometric features
are heavily right-skewed (local density spans $0.03$ to $35$), we
rank-transform features within each fold and report Spearman rank
correlation, which is the relevant metric for a pre-training risk index that
only needs to order datasets.

DPRI ranks datasets by risk with a Spearman correlation of $0.61$ under an
unbiased nested cross-validation that selects the feature subset on training
folds only (Table~\ref{tab:regression}). The fixed-feature LOO estimate is
$0.52$ with all features and $0.61$ with the geometric-core-plus-dimension
subset. Random Forest regression does worse ($\rho=0.36$), overfitting the
small sample. The prediction is moderate, not precise. Coefficients of
determination are not an informative lens here, since the target is a
ranking and the features are rank-transformed.

\begin{table}[t]
\centering
\caption{DPRI ranks $\mathrm{Risk}(D)$ across 31 datasets (Spearman $\rho$,
         leave-one-dataset-out). The nested-CV number is the unbiased
         headline.}
\label{tab:regression}
\small
\begin{tabular}{lc}
\toprule
\textbf{Estimator} & \textbf{Spearman} $\rho$ \\
\midrule
Nested-CV (unbiased, headline)        & $0.61$ \\
Ridge, all six features (fixed)       & $0.52$ \\
Ridge, geometric core + dimension     & $0.61$ \\
Random Forest, all features           & $0.36$ \\
\bottomrule
\end{tabular}
\end{table}

Table~\ref{tab:dpri_features} illustrates the relationship on representative
datasets. High-risk datasets tend to have high sample uniqueness and low
density (Texas100, Purchase100), and low-risk datasets the reverse (Gowalla,
Mushroom). The relationship is a tendency rather than a law: Creditg has
high uniqueness yet only moderate risk, and Mushroom has moderate uniqueness
yet the lowest risk. These exceptions are why the correlation is $0.61$
rather than $1.0$.

\begin{table}[t]
\centering
\caption{DPRI features and $\mathrm{Risk}(D)$ for representative datasets,
         ordered by risk. $u$: sample uniqueness, $\rho$: local density,
         $S$: class separability, $d$: dimensionality.}
\label{tab:dpri_features}
\small
\setlength{\tabcolsep}{4pt}
\begin{tabular}{lrrrrr}
\toprule
\textbf{Dataset} & $\bar{u}$ & $\bar{\rho}$ & $S$ & $d$ & \textbf{Risk} \\
\midrule
Texas100     & $55.7$ & $0.03$ & $-0.25$ & $6169$ & $0.92$ \\
Purchase100  & $28.0$ & $0.04$ & $-0.02$ &  $600$ & $0.85$ \\
Creditg      &  $6.9$ & $0.15$ & $ 0.02$ &   $74$ & $0.77$ \\
Adult        &  $1.0$ & $1.55$ & $-0.01$ &   $14$ & $0.70$ \\
MNIST        & $16.2$ & $0.08$ & $-0.05$ &  $784$ & $0.66$ \\
Magic        &  $0.8$ & $1.49$ & $ 0.12$ &   $10$ & $0.63$ \\
Movielens    &  $0.2$ & $6.10$ & $ 0.12$ &    $5$ & $0.53$ \\
Gowalla      &  $0.1$ & $34.9$ & $ 0.09$ &    $6$ & $0.51$ \\
Mushroom     &  $3.8$ & $0.28$ & $ 0.09$ &  $138$ & $0.51$ \\
\bottomrule
\end{tabular}
\end{table}

\begin{figure}[t]
  \centering
  \dprifig{figures/regression_plot.pdf}{\linewidth}
  \caption{Predicted vs.\ observed $\mathrm{Risk}(D)$ under
           leave-one-dataset-out cross-validation. Each point is one of the
           31 datasets.}
  \label{fig:regression}
\end{figure}

\subsection{Which Factors Carry the Signal}
\label{sec:single_factor}

We assess each feature by its standalone Spearman correlation with
$\mathrm{Risk}(D)$ (Table~\ref{tab:single_factor}). A drop-one ablation is
not meaningful here, because the features are collinear and the linear
baseline is weak, so removing one feature lets a correlated feature absorb
its signal. Standalone rank correlation is robust to both.

Three geometric features are individually significant: local density,
class separability, and sample uniqueness (all $p < 0.005$). Outlier score,
feature entropy, and dimensionality are not individually significant at
$n=31$, though Random Forest assigns them nonzero importance, suggesting
they contribute weak or nonlinear signal. The headline result is that
no single statistic determines risk, and that the significant factors are
precisely the geometric ones our theory predicts.

\begin{table}[t]
\centering
\caption{Standalone ranking power of each DPRI feature (Spearman vs.\
         $\mathrm{Risk}(D)$, $n=31$).}
\label{tab:single_factor}
\small
\begin{tabular}{lcc}
\toprule
\textbf{Feature} & \textbf{Spearman} $\rho$ & \textbf{$p$-value} \\
\midrule
Local density $\bar{\rho}$      & $-0.565$ & $0.001$ \\
Class separability $S$          & $-0.543$ & $0.002$ \\
Sample uniqueness $\bar{u}$     & $+0.506$ & $0.004$ \\
Dimensionality $\log d$         & $+0.294$ & $0.108$ \\
Outlier score $\bar{o}$         & $+0.292$ & $0.111$ \\
Feature entropy $H$             & $-0.084$ & $0.654$ \\
\bottomrule
\end{tabular}
\end{table}

\subsection{Finding 1: Dataset Structure Governs Risk Among Realistic Models}
\label{sec:finding1}

We decompose the variation in MIA AUC by the spread of group-mean AUC across
each factor. Across all nine configurations, dataset choice accounts for
$36\%$ of the variation, model architecture for $39\%$, and the attack
method for $24\%$. Taken at face value, the model appears to rival the
dataset.

This parity is almost entirely an artifact of a single model. The Random
Forest, trained without a depth limit, memorizes its training set and
saturates near AUC $1.0$ on nearly every dataset. Its mean AUC is $0.75$,
against $0.56$ for the MLP and $0.62$ for XGBoost. This is an adversarial
worst case, not a configuration that practitioners deploy. Restricting the
decomposition to the regularized models that are actually used, MLP and
XGBoost, dataset structure rises to $49.5\%$ of the variation, twice the
model share ($24.4\%$) and well ahead of attack ($26.1\%$).

We therefore read Finding~1 as follows. Across realistic models, dataset
structure governs membership inference risk more than the choice of model,
and far more than the choice of attack. A deliberately overfit model can
erase that margin, which is itself a caution. Reporting a single attack
number on a single dataset conflates three distinct sources of variation,
and the dataset is not the incidental one.

\subsection{Finding 2: Benchmarks Are Geometric Outliers}
\label{sec:finding2}

The two most widely used MIA benchmarks, Purchase100 and Texas100, are
geometric outliers among the 31 datasets. Their sample uniqueness ($28.0$
and $55.7$) is the highest of all datasets we study, far above the median
of roughly $1.5$, and their risk ($0.85$ and $0.92$) is likewise the top
two. A model evaluated only on these benchmarks is being attacked in an
unrepresentatively high-risk regime.

We temper this in two ways. First, the effect is not exclusive to
benchmarks: high-dimensional image datasets such as MNIST and Fashion-MNIST
also have elevated uniqueness, because high dimensionality sparsifies the
data. Second, with only two benchmark datasets a formal distributional test
has little power. The claim we make is therefore concrete and limited:
on the single most important geometric axis, sample uniqueness, the two
canonical MIA benchmarks sit at the extreme of the entire 31-dataset range.
Corollary~\ref{cor:benchmark} explains why subsampled, class-balanced
benchmarks are expected to have inflated uniqueness from first principles.

% =============================================================================
\section{Security Implications}
\label{sec:implications}
% =============================================================================

\subsection{Adversarial Use: DPRI as a Reconnaissance Tool}

An adversary with access to multiple candidate datasets---for example,
a researcher evaluating which of several publicly released models to
attack---can compute DPRI in minutes without any model training and
prioritize high-risk targets.

We simulate this scenario.
An adversary has a budget of $B$ MIA queries (each query = training and
attacking one model on one dataset).
Under random dataset selection, the expected AUC return is the
population mean.
Under DPRI-guided selection (attack highest-DPRI first), the adversary
reaches the top-$k$ risk datasets in expectation with
\TODO{X}$\times$ fewer queries (Figure~\ref{fig:adv_query}).

\begin{figure}[t]
  \centering
  \includegraphics[width=0.85\linewidth]{../results/regression/adv_query_efficiency.pdf}
  \caption{Adversarial query efficiency: DPRI-guided vs.\ random dataset
           selection. Shaded region = 95\% CI over 1,000 random orderings.}
  \label{fig:adv_query}
\end{figure}

\paragraph{Mitigation.}
Data holders can compute DPRI before releasing a dataset or a model
trained on it.
A high DPRI score should trigger additional safeguards: stronger DP
guarantees, data minimization, or access controls on the trained model.

\subsection{DP Budget Recommendation}

Based on the empirical relationship between DPRI and DP efficacy
(Finding~3), we propose a simple DPRI-adaptive $\varepsilon$ rule:

\begin{equation}
  \varepsilon_{\mathrm{rec}}(D) =
  \begin{cases}
    \varepsilon_0 & \text{if } \mathrm{DPRI}(D) < 0.35 \\
    \varepsilon_0 / 3 & \text{if } 0.35 \leq \mathrm{DPRI}(D) < 0.65 \\
    \varepsilon_0 / 10 & \text{if } \mathrm{DPRI}(D) \geq 0.65
  \end{cases}
  \label{eq:eps_rec}
\end{equation}

where $\varepsilon_0$ is the practitioner's nominal budget.
This rule is calibrated so that all datasets achieve approximately
equal post-DP MIA AUC, regardless of their native structure.
The thresholds are derived from our empirical calibration on
Table~\ref{tab:dp_calibration} and should be recalibrated as more
datasets are added to the benchmark.

\subsection{Policy Implication}

Current regulatory frameworks for machine learning privacy (e.g., GDPR
\S{}22, HIPAA de-identification) treat privacy as a property of the
algorithm, not the data.
Our findings imply that this is insufficient: two organizations applying
identical DP-SGD with identical $\varepsilon$ achieve systematically
different levels of membership inference protection depending on the
native structure of their data.

We recommend that data governance frameworks adopt DPRI (or a
standardized equivalent) as a required pre-training assessment,
analogous to environmental impact assessments required before
construction projects.
The DPRI tool we release makes this practical at near-zero computational
cost.

\subsection{Extensions}

\paragraph{Federated learning.}
In federated settings, each client's local dataset has its own DPRI.
Clients with high DPRI should contribute more noise in the local DP
step, or be excluded from the federation if their risk exceeds a
threshold.

\paragraph{Generative models.}
DPRI is defined on the training data distribution and is
model-agnostic; it applies equally to GANs, diffusion models, and LLMs.
We leave empirical validation on generative model training data as
future work.

% =============================================================================
\section{Related Work}
\label{sec:related}
% =============================================================================

\paragraph{Membership inference attacks.}
Shokri et al.~\cite{shokri2017membership} introduced the shadow-model
attack and established Purchase100 and Texas100 as standard benchmarks.
Subsequent work strengthened attacks via label-only
access~\cite{choquette2021label}, enhanced
thresholds~\cite{ye2022enhanced}, and the likelihood-ratio
formulation~\cite{carlini2022lira} that we use as ground truth.
Salem et al.~\cite{salem2019ml} showed that attacks are robust to
imperfect shadow model knowledge.
Our work does not propose a new attack. Rather, it ranks datasets by their
attack success from data properties alone, before any model is trained.

\paragraph{Privacy risk theory.}
Yeom et al.~\cite{yeom2018privacy} proved that MIA advantage is bounded
by the generalization gap, linking privacy to overfitting.
Sablayrolles et al.~\cite{sablayrolles2019whitebox} characterized the
Bayes-optimal MIA as a likelihood ratio test, which forms the basis of
our Theorem~\ref{thm:main}.
Our contribution extends these results by showing that the likelihood
ratio, and hence the MIA advantage, is bounded, under explicit modeling
assumptions, by pre-training distributional statistics, not only by
post-training model properties.

\paragraph{Memorization and long-tail data.}
Feldman~\cite{feldman2020memorization} proved that learning algorithms
must memorize atypical (long-tail) training examples to achieve near-optimal
generalization, and connected this to per-sample privacy risk.
Carlini et al.~\cite{carlini2021extracting} showed that LLMs memorize
rare training examples verbatim.
Our work differs along three dimensions.
Unit of analysis: we study dataset-level statistics rather than
per-sample certificates.
Timing: we operate before training, with no trained model.
Output: a dataset-level risk score used to rank datasets, not a post-hoc
per-sample memorization label.
Our results are consistent with Feldman's. The datasets we find most at
risk are those with high sample uniqueness, the aggregate counterpart of
atypical, memorization-prone examples. We complement that line of work by
asking, at the dataset level and before training, how far this geometry
predicts risk, and we find that it does so moderately but significantly.

\paragraph{Differential privacy.}
DP-SGD~\cite{abadi2016deep} is the standard mechanism for training
private neural networks.
Jagielski et al.~\cite{jagielski2020auditing} show empirically that the
privacy guarantees of DP-SGD are tight.
Zanella-Béguelin et al.~\cite{zanella2023bayesian} propose Bayesian
estimation of the actual $\varepsilon$ achieved.
None of these works ask how the benefit of a given $\varepsilon$
varies with dataset structure. We raise this as a direction motivated by our
bound (Section~\ref{sec:implications}), though we do not evaluate DP-SGD in
this work.

\paragraph{Dataset auditing and data valuation.}
Long et al.~\cite{long2020pragmatic} propose a pragmatic approach to
membership inference that accounts for dataset characteristics.
Our work goes further by predicting risk before training and providing a
formal theoretical basis through Theorem~\ref{thm:main}.

% =============================================================================
\section{Conclusion}
\label{sec:conclusion}
% =============================================================================

We introduced DPRI, a pre-training dataset privacy risk estimator grounded
in a formal bound on membership inference advantage.
Across 31 datasets in five domains, DPRI ranks MIA vulnerability with an
unbiased cross-validated Spearman correlation of $0.61$, using
distributional features that require no model training.
The prediction is moderate rather than precise, which is itself a finding.
Privacy risk is significantly shaped by data geometry, yet no single
pre-training statistic determines it.
Two further findings challenge established practice. First, the benchmarks
that anchor a decade of membership inference research are geometric outliers
that do not represent typical deployed data. Second, among the regularized
models practitioners actually deploy, dataset structure governs risk more
than model architecture and far more than the choice of attack, so reporting
an attack number on a single benchmark conflates three distinct sources of
variation, of which the dataset is the largest.

These results suggest a shift in how the community measures and reports
membership inference risk.
We call for benchmark diversity in MIA evaluations and for dataset-level
risk assessment as a standard pre-training step, while cautioning that
pre-training prediction is informative but not a substitute for direct
measurement.

The DPRI tool is provided as an anonymized artifact and will be released
as open source upon publication.


\bibliographystyle{IEEEtran}
\bibliography{references}

\appendix
% =============================================================================
\section{Proof of Theorem~\ref{thm:main}}
\label{app:proofs}
% =============================================================================

We restate and prove Theorem~\ref{thm:main} in full.

\begin{theorem}[DPRI Upper Bound on MIA Advantage, restated]
Let $f_\theta$ be trained on $D$ by minimizing empirical risk with an
$L$-Lipschitz loss $\ell$.
Let $P_{\mathrm{in}}$ and $P_{\mathrm{out}}$ be the distributions of
$\ell(f_\theta, x)$ for members and non-members, respectively.
Assume $P_{\mathrm{out}}$ is sub-Gaussian with variance proxy $\sigma^2$.
Then for any sample $x \in D$:
\begin{equation}
  \mathrm{Adv}(\mathcal{A}^*, f_\theta, x)
  \leq
  \frac{L \cdot u(x)}{\sigma \sqrt{2} \cdot \rho(x)^{1/d}},
  \label{eq:bound_app}
\end{equation}
where $u(x) = \min_{x' \neq x} \|x - x'\|_2$ and
$\rho(x)$ is the $k$-NN density estimate of
Equation~\eqref{eq:density}. In the DPRI implementation we use the
monotone surrogate $1/r_k(x)$, which has the same rate under (A3) and
avoids evaluating high-dimensional ball volumes, which underflow for
large $d$.
\end{theorem}

\noindent\textbf{Assumptions.} We make three modeling assumptions. (A1) The
member and non-member loss distributions $P_{\mathrm{in}}$ and
$P_{\mathrm{out}}$ are Gaussian with a common variance $\sigma^2$ and means
$\mu_{\mathrm{in}}, \mu_{\mathrm{out}}$. (A2) The loss $\ell(f_\theta,\cdot)$
is $L$-Lipschitz in its input. (A3) The local nearest-neighbour distance and
density satisfy $r_k(x) \asymp \rho(x)^{-1/d}$, which holds when the density
is approximately constant over the neighbourhood~\cite{biau2015lectures}.

\begin{proof}
\textbf{Step 1 (advantage to KL).}
For the optimal test between $P_{\mathrm{in}}$ and $P_{\mathrm{out}}$,
the advantage equals the total variation distance, which Pinsker's
inequality~\cite{tsybakov2009introduction} bounds by the KL divergence:
\begin{equation}
  \mathrm{Adv}(\mathcal{A}^*) \leq
  \sqrt{\tfrac{1}{2}\,\mathrm{KL}(P_{\mathrm{in}} \| P_{\mathrm{out}})}.
  \label{eq:bh_app}
\end{equation}

\textbf{Step 2 (KL via the loss gap).}
Under (A1), the KL divergence between two Gaussians of common variance is
exact:
\begin{equation}
  \mathrm{KL}(P_{\mathrm{in}} \| P_{\mathrm{out}})
  = \frac{(\mu_{\mathrm{in}} - \mu_{\mathrm{out}})^2}{2\sigma^2}.
\end{equation}
By (A2), the mean losses of a member $x$ and its nearest out-of-sample point
$x^* = \arg\min_{x' \notin D}\|x-x'\|_2$ differ by at most $L\|x-x^*\|_2$.
Therefore
\begin{equation}
  \mathrm{KL}(P_{\mathrm{in}} \| P_{\mathrm{out}})
  \leq \frac{L^2 \|x - x^*\|_2^2}{2\sigma^2}.
  \label{eq:kl_app}
\end{equation}

\textbf{Step 3 (distance via density).}
The member-to-nearest-neighbour distance is $u(x)$ by definition, and by (A3)
the spacing to a fresh draw from $\mathcal{P}$ scales as $\rho(x)^{-1/d}$.
Combining these into the geometric factor $u(x)/\rho(x)^{1/d}$, substituting
into~\eqref{eq:kl_app}, and applying~\eqref{eq:bh_app},
\begin{equation}
  \mathrm{Adv}(\mathcal{A}^*) \leq
  \sqrt{\frac{L^2 u(x)^2}{2\sigma^2 \rho(x)^{2/d}}}
  = \frac{L\, u(x)}{\sigma\sqrt{2}\,\rho(x)^{1/d}}.
\end{equation}
Averaging over $x \in D$ gives the dataset-level bound~\eqref{eq:risk_bound}.
\end{proof}

\noindent\textbf{Remark.} The bound is a modeling result rather than a
worst-case guarantee. It rests on the Gaussian score model (A1) and the local
density-distance relation (A3), both standard but approximate. Its value is
the qualitative prediction it makes, that MIA advantage grows with sample
uniqueness and falls with local density, which Section~\ref{sec:results}
tests directly against measured AUC and finds to hold as a moderate,
statistically significant trend.

\begin{proposition}[Tightness]
\label{prop:tight}
The bound~\eqref{eq:bound_app} is tight up to constants for datasets
drawn from a uniform distribution over a $d$-dimensional hypercube.
\end{proposition}
\begin{proof}
For uniform data, $\rho(x) = c$ (constant) and $u(x) \sim n^{-1/d}$
by standard spacing theory.
The Bayes-optimal MIA advantage for a uniformly distributed dataset
also scales as $O(n^{-1/d})$ (from the Yatracos class bound).
Hence the bound and the true advantage have the same rate.
\end{proof}

Section~\ref{sec:tightness_main} validates
Proposition~\ref{prop:tight} empirically on uniform hypercube data:
the bound factor's fitted decay slopes match the predicted $-1/d$
within $0.03$ for $d \in \{2,4,8\}$.

% =============================================================================
\section{Dataset Details}
\label{app:datasets}
% =============================================================================

The 31 datasets span medical (e.g.\ Breast-W, Pima Diabetes), image (MNIST,
Fashion-MNIST, Digits, OptDigits, PenDigits), network (KDDCup), forest
cover (CovType), credit (German Credit), and a range of tabular sources
(OpenML), together with the canonical MIA benchmarks Purchase100 and
Texas100~\cite{shokri2017membership} and the mobility and recommendation
datasets Gowalla~\cite{cho2011friendship} and
MovieLens~\cite{movielens}.
Datasets larger than 30k samples are subsampled to 30k to keep the
attack-model grid tractable. All datasets are publicly available and loaded
through scikit-learn, OpenML, or their original public sources.


\end{document}
